
\chapter{Diseño}
\label{cap:diseno}

En este capítulo se describe el diseño de la solución. También muestra cómo incluir código fuente y ecuaciones.

\section{Código fuente}

El código se incluye con el entorno \texttt{lstlisting}. Los estilos disponibles (\texttt{C}, \texttt{C++}, \texttt{Java}, \texttt{Python}, \texttt{PHP}, \texttt{R}, \texttt{Matlab}, \texttt{Latex}, \texttt{Consola}\ldots) están en \texttt{config/listing-styles.sty}. El \cref{lst:python} es un ejemplo.

\begin{lstlisting}[style=Python, caption={Ejemplo de código en Python.}, label={lst:python}]
def factorial(n: int) -> int:
    """Calcula el factorial de n de forma recursiva."""
    if n <= 1:
        return 1
    return n * factorial(n - 1)

print(factorial(5))  # Imprime 120
\end{lstlisting}

También se puede incluir un fichero completo:

\begin{lstlisting}[style=Latex]
\lstinputlisting[style=Java]{codigo/Main.java}
\end{lstlisting}

\section{Ecuaciones}

Las ecuaciones numeradas se escriben con el entorno \texttt{equation} y se referencian igual que el resto de elementos, como la \cref{eq:complejidad}:

\begin{equation}
  T(n) = \sum_{i=1}^{n} i = \frac{n(n+1)}{2} \in O(n^2)
  \label{eq:complejidad}
\end{equation}
